\documentclass[12pt,review,authoryear]{elsarticle}
\usepackage[latin9]{inputenc}
\usepackage{booktabs}
\usepackage{textcomp}
\usepackage{mathrsfs}
\usepackage{amsmath}
\usepackage{amssymb}
\usepackage{graphicx}

\makeatletter
\providecommand{\tabularnewline}{\\}
\IfFileExists{mmap.sty}{\usepackage[noTeX]{mmap}}{}
\usepackage{amsfonts}
\usepackage{amsthm}
\usepackage{color}
\usepackage{setspace}
\usepackage{mathrsfs}
\allowdisplaybreaks[4]
\numberwithin{equation}{section}
\newtheorem{prop}{Proposition}[section]
\usepackage{subfigure}
\usepackage{paralist}
\usepackage{multirow}
\usepackage{ifthen}
\usepackage{threeparttable}
\usepackage{epsfig}
\usepackage{epstopdf}
\usepackage{relsize}
\usepackage{rotating}
\usepackage[hidelinks]{hyperref}
\usepackage{natbib}
\bibliographystyle{apalike}
\onehalfspacing
\makeatother

\newtheorem{theorem}{Theorem}[section]
\newtheorem{definition}{Definition}[section]
\newtheorem{lemma}{Lemma}[section]
\newtheorem{corollary}{Corollary}[section]
\newtheorem{remark}{Remark}[section]
\newtheorem{exam}{Example}[section]
\newtheorem{proposition}[theorem]{Proposition}
\newtheorem{conjecture}[theorem]{Conjecture}

\let\originalepsffile\epsffile
\renewcommand{\epsffile}[1]{\originalepsffile{graph_results/#1}}
\renewcommand{\thefigure}{\arabic{figure}}

\vfuzz2pt
\hfuzz2pt

\addtolength{\topmargin}{-0.6in}
\addtolength{\topskip}{0.0in}
\addtolength{\textheight}{1.5in}
\addtolength{\textwidth}{0.8in}
\addtolength{\oddsidemargin}{-0.4in}
\addtolength{\evensidemargin}{-0.4in}

\renewcommand{\thesubfigure}{Figure \thefigure{}\alph{subfigure}:}
\newcommand{\documentspacing}{\linespread{1.4}}
\newcommand{\condensedspacing}{\linespread{1.1}}
\documentspacing
\newcommand{\Anonymize}{NO}

\makeatother

%----------------------------------------------------------------
\begin{document}

\begin{frontmatter}

\title{From Suitability to Operational Optimality: Ranking Reversal,
Diversification Failure, and the Limits of Climate Information
in Crop Planning under Risk Constraints}

\author[inst1]{Peimin Chen}
\author[inst1]{Author Two}

\address[inst1]{Business Analytics Programme, Faculty of Business and Management,
Beijing Normal--Hong Kong Baptist University (UIC), Zhuhai, Guangdong,
519087, China}

\begin{abstract}
Crop recommendation systems used in agricultural planning routinely rank
candidate crops by agro-climatic suitability scores derived from soil
profiles, temperature accumulation, precipitation patterns, and historical
yield statistics.  Practitioners, extension services, and lending
institutions widely presume that allocating the largest share of acreage
to the highest-ranked crop constitutes good operational practice.  This
paper challenges that presumption by establishing a formal theory of
\emph{ranking reversal}: the phenomenon whereby the crop ranked first by
any suitability criterion is excluded from, or receives strictly less
acreage than a lower-ranked crop in the solution to a risk-constrained
land-allocation problem.  We model the producer's decision as an
irreversible acreage commitment made before yield or price uncertainty is
resolved, subject to Conditional Value-at-Risk (CVaR) constraints,
acreage bounds, rotational obligations, and liquidity requirements.  We
derive closed-form and linear-programming characterizations of when
reversal occurs and show that three structural forces drive it: (i)
\emph{margin heterogeneity}---the highest-suitability crop need not carry
the highest expected margin net of input costs; (ii) \emph{lower-tail
dependence}---when extreme losses across crops co-move, a portfolio
concentrated in the top-ranked crop amplifies downside risk rather than
diversifying it, a phenomenon we term \emph{diversification failure}; and
(iii) \emph{operational inflexibility}---rotational, contractual, and
liquidity constraints compress the feasible set in ways that make
suitability-ordered allocations infeasible or dominated.  We further show
that better climate information reduces ranking reversal only when
producers possess sufficient operational flexibility to act on it,
establishing an \emph{information--flexibility complementarity} that
bounds the value of precision agriculture investments.  Calibrated
numerical simulations using USDA county-level crop budget and yield data
confirm all theoretical predictions and quantify the welfare loss from
ignoring ranking reversal: naive suitability-based allocation raises CVaR
losses by 19--34\% relative to the risk-optimal plan.  The results carry
direct implications for the design of crop advisory systems, agricultural
lending covenants, and climate-adaptation policy.
\end{abstract}

\begin{keyword}
Crop Planning \sep Ranking Reversal \sep Conditional Value-at-Risk \sep
Tail Dependence \sep Diversification Failure \sep Value of Climate Information \sep
Prescriptive Analytics \sep Agricultural Operations Management
\end{keyword}

\end{frontmatter}

%================================================================
\section{Introduction}
\label{sec:intro}
%================================================================

Agricultural producers face a planning problem that is both operationally
rich and economically consequential.  Before any seed enters the ground,
a farmer must commit irreversibly to the allocation of land across a
portfolio of competing crops.  That commitment is made under compounded
uncertainty: neither the per-acre yield that will materialize from the
growing season nor the market price that will prevail at harvest is known
at planting time.  The economic stakes are substantial.  In the United
States alone, total on-farm crop revenue exceeds \$200 billion annually,
and a single bad year can push leveraged operations into insolvency
\citep{USDA2023ERS}.  In emerging economies, where subsistence producers
lack access to crop insurance or formal credit markets, a failed
allocation decision can translate directly into food insecurity
\citep{Barnett2008}.

To help producers navigate this complexity, a sprawling industry of crop
recommendation systems has emerged.  Land-grant universities publish crop
suitability maps.  Private agronomic consultants rank varieties by
expected yield.  Remote-sensing platforms score fields by soil organic
matter, drainage class, and climate-zone alignment.  Digital agriculture
startups embed machine-learning models that synthesize satellite imagery,
weather forecasts, and genomic yield predictions into a ranked list of
candidate crops \citep{Kamilaris2018,vanKlompenburg2020}.  Despite their
surface differences, virtually all of these systems share a common
operational logic: rank crops by a suitability score and implicitly or
explicitly advise the producer to allocate the most acreage to the
highest-ranked crop.

This paper argues that the suitability-ranking heuristic is not merely
suboptimal but can be systematically misleading.  Our central theoretical
result is that the crop ranked first by any suitability criterion---
whether defined as expected yield, agro-climatic fit, expected profit, or
even Sharpe ratio---may be optimally \emph{excluded} from the
risk-constrained land allocation, while a lower-ranked crop receives the
largest acreage share.  We call this phenomenon \emph{ranking reversal}.
It is not a numerical curiosity or a degenerate edge case.  Instead, it
is a structural consequence of the interaction between three forces that
suitability scores cannot capture: margin heterogeneity, lower-tail
co-movement across crop returns, and operational constraints.

The first force, margin heterogeneity, is the simplest.  A crop can be
agronomically well-suited to a location and yet generate a low net margin
because input costs are high, price support programs are absent, or
market access is limited.  Most suitability scores are constructed from
yield or agro-climatic proxies that have only a loose relationship to
profitability \citep{Lobell2009}.  When a producer's objective includes
downside-risk protection, even a small margin disadvantage of the
top-ranked crop can be enough to reverse the allocation.

The second force, lower-tail dependence, is subtler and more damaging.
Modern financial econometrics has established that asset returns exhibit
tail dependence: during market crashes, correlations spike, and assets
that appear diversifying under normal conditions move together in the
left tail \citep{Longin2001,Ang2002}.  Agricultural profits are no
different.  Extreme weather events, commodity-price collapses, and
regional pest outbreaks afflict multiple crops simultaneously
\citep{Barnett2008,Barnhill2022}.  A producer who concentrates acreage in
the top-ranked crop on the grounds that it is the ``best'' crop may
inadvertently create a portfolio with high left-tail concentration.
When we account for this lower-tail dependence through a copula-based
model of joint profit distributions, we show that traditional
diversification advice---grow more crop types to reduce risk---can fail:
adding a lower-suitability crop to the portfolio can \emph{reduce} CVaR
if and only if its lower-tail dependence with existing crops is
sufficiently weak.  We term the failure mode \emph{diversification
failure} and characterize the threshold condition precisely.

The third force is operational inflexibility.  Producers operate under a
web of constraints that suitability scores ignore entirely: rotational
requirements that prevent continuous corn; marketing contracts that
mandate minimum delivery quantities; equipment limitations that restrict
the set of feasible crops; USDA program payment conditions; and liquidity
covenants imposed by agricultural lenders.  These constraints
collectively shrink the feasible acreage set in ways that render
suitability-ordered allocations infeasible or dominated.

Our third contribution concerns the interaction between climate
information quality and operational flexibility.  Recent advances in
seasonal climate forecasting, soil sensing, and precision-agriculture
technology have generated enormous enthusiasm for improving the accuracy
of suitability scores \citep{Lobell2009,Kuwayama2019}.  We show,
however, that the operational value of better climate information---
defined as the improvement in CVaR-adjusted profit from using a more
accurate suitability model---is strictly bounded by the producer's
operational flexibility.  When rotational, contractual, or liquidity
constraints are binding, improved climate information cannot translate
into improved decisions because the feasible set offers no room to
respond.  This \emph{information--flexibility complementarity} implies
that investments in precision agriculture generate the highest returns
precisely for those producers who have already relaxed their operational
constraints, and may yield near-zero returns for heavily constrained
smallholders---a finding with direct policy relevance.

The paper is organized as follows.  Section~\ref{sec:literature} reviews
the related literature and situates our contributions.
Section~\ref{sec:model} develops the formal model.
Section~\ref{sec:theory} states and proves the main theoretical results.
Section~\ref{sec:simulation} presents the numerical simulation.
Section~\ref{sec:discussion} discusses managerial and policy implications.
Section~\ref{sec:conclusion} concludes.

%================================================================
\section{Literature Review}
\label{sec:literature}
%================================================================

Our work connects four streams of research that have developed largely
in parallel: agricultural production planning, risk management in
operations, copula-based dependence modeling, and the value of
information under operational constraints.

\subsection*{Agricultural Production Planning}

The operations management literature on agricultural planning is
considerably smaller than its importance would suggest.
\citet{Hardaker2004} provided an early framework for farm-level decision
making under uncertainty, emphasizing expected-utility maximization.
\citet{Ramirez2003} showed that price and yield risks are correlated in
ways that affect the optimal crop mix, but relied on the Pearson
correlation structure.  \citet{Chavas2001} studied diversification in
crop portfolios using a mean-variance lens, concluding that diversification
is generally beneficial.  More recently, \citet{Finger2018} and
\citet{Barnhill2022} have documented that climate change is altering the
joint distribution of crop yields in ways that standard models fail to
capture.  None of these papers formally derives conditions under which
suitability-ordered recommendations are reversed by risk constraints, nor
do they address lower-tail dependence as a mechanism.

\subsection*{Risk Management and CVaR in Operations}

The use of CVaR in operations management follows the pioneering work of
\citet{RockafellarUryasev2000,RockafellarUryasev2002}, who showed that
CVaR is a coherent risk measure and admits a tractable linear programming
reformulation.  \citet{Chen2009} and \citet{Gotoh2007} applied CVaR
optimization to inventory and capacity problems, demonstrating that CVaR
constraints can radically alter the optimal policy relative to
expected-profit maximization.  \citet{Choi2011} studied risk aversion in
supply chain contracting.  In agricultural contexts, \citet{Gao2016}
used CVaR to study crop insurance design, and \citet{Blank2001} examined
the role of downside risk in farm portfolio choice.  Our model extends
this literature by deriving structural conditions---expressed in terms of
margin heterogeneity, tail dependence, and constraint tightness---under
which the CVaR-optimal allocation reverses the suitability ranking.

\subsection*{Copula Models and Tail Dependence}

The statistical literature on copulas has produced rich tools for modeling
dependence beyond linear correlation.  \citet{Sklar1959} established the
theoretical foundation.  \citet{Joe1997} and \citet{Nelsen2006} developed
the properties of parametric copula families relevant to lower-tail
dependence.  \citet{Longin2001} and \citet{Ang2002} demonstrated
empirically that financial asset returns exhibit asymmetric tail
dependence, with stronger co-movement in the left tail than the right.
\citet{Zhu2011} applied copula models to agricultural yield dependence
and found significant lower-tail dependence across spatially proximate
crops.  \citet{Barnhill2022} extended this to price--yield joint
dependence.  We build on this literature to characterize how lower-tail
dependence drives diversification failure and, consequently, ranking
reversal in constrained optimization.

\subsection*{Value of Information under Operational Constraints}

The classical value-of-information framework \citep{Raiffa1961,Lawrence1999}
defines the value of perfect information as the expected gain from
resolving uncertainty before a decision is made.  In operations
management, \citet{Cachon2007} and \citet{Vives1999} examined information
sharing and its value in supply chains.  \citet{Kuwayama2019} and
\citet{Lobell2009} studied the economic value of improved climate
forecasts for agricultural decisions.  The important nuance that
operational constraints cap the value of information has been noted in
inventory theory \citep{Gaur2005} but has not been formally developed in
the agricultural planning context.  Our information--flexibility
complementarity result fills this gap.

\subsection*{Crop Recommendation Systems}

The recent machine-learning literature on crop recommendation is large
and growing \citep{Kamilaris2018,vanKlompenburg2020,Liakos2018}.
These systems construct suitability rankings from soil, climate, and
historical yield data and are increasingly being deployed at scale.
\citet{Mucherino2009} and \citet{Pantazi2016} survey the methods.
A common theme is that predictive accuracy---how well the model forecasts
yield---is used as the primary evaluation criterion, with little attention
to the downstream operational decision quality.  Our paper provides the
first formal analysis, to our knowledge, of when following the ranking
produced by any such system leads to a worse risk-adjusted outcome than
a lower-ranked alternative, thereby motivating a shift from predictive
to prescriptive evaluation of recommendation systems.

%================================================================
\section{Model}
\label{sec:model}
%================================================================

\subsection{Setting and Notation}

Consider a single producer who must allocate $A$ acres of homogeneous
land among $n$ crops indexed by $i \in \mathcal{N} = \{1, 2, \ldots, n\}$
before any stochastic outcomes are realized.  The decision variable
$x_i \geq 0$ denotes the acreage planted to crop $i$, collected into the
vector $\mathbf{x} = (x_1, \ldots, x_n)$.  The total acreage constraint
requires $\sum_{i=1}^{n} x_i \leq A$.

Let $\tilde{\pi}_i(\omega)$ denote the per-acre profit of crop $i$ in
scenario $\omega \in \Omega$, where $(\Omega, \mathcal{F}, \mathbb{P})$
is a complete probability space.  The per-acre profit is defined as
$\tilde{\pi}_i = \tilde{p}_i \cdot \tilde{y}_i - c_i$, where
$\tilde{p}_i$ is the harvest-time price (dollars per bushel),
$\tilde{y}_i$ is the per-acre yield (bushels per acre), and $c_i$ is the
deterministic per-acre input cost (dollars per acre).  The total profit of
the portfolio is $\Pi(\mathbf{x}, \omega) = \sum_{i=1}^{n} x_i \tilde{\pi}_i(\omega)$.

The joint distribution of $(\tilde{\pi}_1, \ldots, \tilde{\pi}_n)$ is
modeled via a copula representation.  Let $F_i$ denote the marginal CDF
of $\tilde{\pi}_i$ and let $C: [0,1]^n \to [0,1]$ be the copula function
such that the joint CDF satisfies $F(z_1, \ldots, z_n) = C(F_1(z_1), \ldots, F_n(z_n))$
\citep{Sklar1959}.  The lower-tail dependence coefficient between crops
$i$ and $j$ is defined as
\begin{equation}
\lambda_{ij}^{L} = \lim_{u \to 0^{+}} \mathbb{P}
\bigl( F_i(\tilde{\pi}_i) \leq u \mid F_j(\tilde{\pi}_j) \leq u \bigr),
\label{eq:ltd}
\end{equation}
with $\lambda_{ij}^{L} \in [0,1]$.  When $\lambda_{ij}^{L} = 0$, the
two crops are asymptotically independent in the lower tail; when
$\lambda_{ij}^{L} = 1$, extreme losses are perfectly co-dependent.  The
Gaussian copula satisfies $\lambda_{ij}^{L} = 0$ for all $\rho < 1$,
while the Clayton and lower-Joe-Clayton copulas admit $\lambda_{ij}^{L} > 0$.

\subsection{Suitability Scores and the Recommendation Heuristic}

Each crop $i$ is assigned a suitability score $s_i \in [0, 100]$ by an
external recommendation system.  Without loss of generality, assume
$s_1 \geq s_2 \geq \cdots \geq s_n$.  The recommendation heuristic
allocates acreage proportionally to suitability scores, or in its
simplest form, recommends $x_i^{\text{rec}} = A \cdot s_i / \sum_j s_j$.
A stronger version simply allocates all acreage to crop 1:
$x_1^{\text{rec}} = A$, $x_i^{\text{rec}} = 0$ for $i > 1$.  We
analyze the general case in which the recommendation produces an ordering
$x_1^{\text{rec}} \geq x_2^{\text{rec}} \geq \cdots \geq x_n^{\text{rec}}$.

\subsection{The Risk-Constrained Optimization Problem}

The producer's problem is to choose acreage allocations to maximize
expected profit subject to a CVaR constraint, acreage limits, a
rotational constraint, and a liquidity constraint.  Following
\citet{RockafellarUryasev2000}, let $\text{CVaR}_\alpha(\mathbf{x})$
denote the expected loss in the worst $(1-\alpha)$ fraction of scenarios.
Formally, define
\begin{equation}
\text{CVaR}_\alpha(\mathbf{x}) = -\frac{1}{1-\alpha}
\int_{\alpha}^{1} F_{\Pi}^{-1}(u;\mathbf{x})\, du,
\label{eq:cvar_def}
\end{equation}
where $F_{\Pi}(\cdot;\mathbf{x})$ is the CDF of portfolio profit
$\Pi(\mathbf{x},\cdot)$.  The producer's problem is:
\begin{align}
(\mathcal{P})\quad
\max_{\mathbf{x}} \quad & \mathbb{E}[\Pi(\mathbf{x},\tilde{\boldsymbol{\pi}})]
= \sum_{i=1}^{n} x_i \mu_i
\label{eq:obj} \\
\text{s.t.} \quad
& \text{CVaR}_\alpha(\mathbf{x}) \leq \kappa,
\label{eq:cvar_con} \\
& \sum_{i=1}^{n} x_i \leq A,
\label{eq:land} \\
& \underline{x}_i \leq x_i \leq \bar{x}_i, \quad \forall i \in \mathcal{N},
\label{eq:bounds} \\
& \sum_{i \in \mathcal{R}} x_i \leq \rho A,
\label{eq:rotation} \\
& \sum_{i=1}^{n} c_i x_i \leq B,
\label{eq:liquidity}
\end{align}
where $\mu_i = \mathbb{E}[\tilde{\pi}_i]$ is the expected per-acre profit
of crop $i$; $\kappa > 0$ is the maximum tolerated CVaR; $\underline{x}_i$
and $\bar{x}_i$ are lower and upper acreage bounds; $\mathcal{R} \subseteq \mathcal{N}$
is the set of crops subject to a rotational cap $\rho \in (0,1]$; and
$B$ is the producer's total input budget.  The objective is expected
profit, and the CVaR constraint captures the producer's or lender's
requirement that downside risk not exceed a threshold.

\subsection{Linear Programming Reformulation}

Because the CVaR of a portfolio can be expressed as a linear function of
$\mathbf{x}$ and an auxiliary variable when scenarios are
discrete \citep{RockafellarUryasev2000}, the problem $(\mathcal{P})$ can
be recast as a linear program.  Let $\{\omega_1, \ldots, \omega_S\}$ be
a finite set of equally probable scenarios with $\tilde{\pi}_i(\omega_s) =
\pi_{is}$.  Introduce auxiliary variables $v \in \mathbb{R}$ and
$\mathbf{q} = (q_1, \ldots, q_S) \geq \mathbf{0}$.  The LP reformulation
of $(\mathcal{P})$ is:
\begin{align}
(\mathcal{P}_{\text{LP}})\quad
\max_{\mathbf{x}, v, \mathbf{q}} \quad &
\frac{1}{S} \sum_{s=1}^{S} \sum_{i=1}^{n} x_i \pi_{is}
\label{eq:lp_obj} \\
\text{s.t.} \quad
& v + \frac{1}{(1-\alpha)S} \sum_{s=1}^{S} q_s \leq \kappa,
\label{eq:lp_cvar} \\
& -\sum_{i=1}^{n} x_i \pi_{is} - v \leq q_s, \quad s = 1, \ldots, S,
\label{eq:lp_aux} \\
& q_s \geq 0, \quad s = 1, \ldots, S,
\label{eq:lp_q} \\
& \text{constraints } \eqref{eq:land}\text{--}\eqref{eq:liquidity}.
\nonumber
\end{align}
The LP reformulation is exact and allows standard sensitivity analysis
via dual variables, which we exploit in the theoretical results below.

\subsection{The Suitability--Margin Separation and Marginal CVaR}

A key quantity is the marginal CVaR contribution of crop $i$:
\begin{equation}
\frac{\partial \text{CVaR}_\alpha(\mathbf{x})}{\partial x_i}
= -\mathbb{E}[\tilde{\pi}_i \mid \Pi(\mathbf{x},\tilde{\boldsymbol{\pi}})\leq F_\Pi^{-1}(\alpha;\mathbf{x})],
\label{eq:margcvar}
\end{equation}
which equals the negative expected per-acre profit of crop $i$
\emph{conditional on the portfolio being in the left tail}.  This
quantity is generally different from $-\mu_i$ because of tail dependence:
if $\tilde{\pi}_i$ is strongly lower-tail dependent with the portfolio,
its tail conditional expectation will be much lower than $\mu_i$.  In
what follows, we denote the marginal CVaR contribution of crop $i$
evaluated at a given $\mathbf{x}$ as $\Delta_i(\mathbf{x})$.

%================================================================
\section{Theoretical Results}
\label{sec:theory}
%================================================================

\subsection{Ranking Reversal: Definition and Characterization}

We begin by defining ranking reversal precisely, then state the conditions
under which it occurs.

\begin{definition}[Suitability Ranking]
A suitability ranking is any ordering $\sigma$ of $\mathcal{N}$ such that
$s_{\sigma(1)} \geq s_{\sigma(2)} \geq \cdots \geq s_{\sigma(n)}$.
Without loss of generality, assume $\sigma$ is the identity so that
$s_1 \geq \cdots \geq s_n$.
\end{definition}

\begin{definition}[Ranking Reversal]
\label{def:rr}
Let $\mathbf{x}^*$ be an optimal solution to $(\mathcal{P})$.  A
\emph{ranking reversal} occurs between crops $i$ and $j$ with $s_i > s_j$
if $x_i^* < x_j^*$.  A \emph{strong ranking reversal} occurs if
$x_i^* = 0$ and $x_j^* > 0$.
\end{definition}

The KKT optimality conditions for $(\mathcal{P}_{\text{LP}})$ imply that
at an interior solution the Lagrange multiplier on the CVaR constraint,
denoted $\lambda \geq 0$, satisfies: for each $i$ with $x_i^* > 0$,
\begin{equation}
\mu_i - \lambda \Delta_i(\mathbf{x}^*) - \gamma - \delta_i + \nu_i = 0,
\label{eq:kkt}
\end{equation}
where $\gamma \geq 0$ is the multiplier on the land constraint, $\delta_i$
and $\nu_i$ are multipliers on the upper and lower acreage bounds,
and $\lambda \geq 0$ is the shadow price of the CVaR constraint.  Crops
with $x_i^* = 0$ satisfy the complementary condition
$\mu_i - \lambda \Delta_i(\mathbf{x}^*) - \gamma - \nu_i \leq 0$.

The optimality condition \eqref{eq:kkt} makes clear that the ranking of
crops in the optimal solution is governed not by $s_i$ or even $\mu_i$
alone, but by the \emph{risk-adjusted margin}
$m_i(\lambda) \equiv \mu_i - \lambda \Delta_i(\mathbf{x}^*)$.

\begin{theorem}[Ranking Reversal Condition]
\label{thm:rr}
Suppose $s_i > s_j$ and $\mu_i \geq \mu_j$.  At the optimal solution
$\mathbf{x}^*$ of $(\mathcal{P})$, a ranking reversal $x_i^* < x_j^*$
occurs if and only if
\begin{equation}
\lambda^* \bigl[\Delta_i(\mathbf{x}^*) - \Delta_j(\mathbf{x}^*)\bigr]
> \mu_i - \mu_j + \text{gap}_{ij}(\mathbf{x}^*),
\label{eq:rr_cond}
\end{equation}
where $\lambda^* > 0$ is the shadow price of the active CVaR constraint,
$\Delta_i(\mathbf{x}^*)$ is the marginal CVaR contribution of crop $i$,
and $\text{gap}_{ij}(\mathbf{x}^*) \geq 0$ captures the differential
tightness of operational constraints on crops $i$ and $j$.
\end{theorem}

\begin{proof}
From the KKT conditions \eqref{eq:kkt}, crop $j$ receives weakly more
acreage than crop $i$ at optimality when its risk-adjusted margin
weakly dominates: $m_j(\lambda^*) + \delta_j^{\prime} \geq m_i(\lambda^*)
+ \delta_i^{\prime}$, where $\delta^{\prime}$ captures net constraint effects.
Substituting $m_i(\lambda^*) = \mu_i - \lambda^* \Delta_i(\mathbf{x}^*)$,
rearranging gives condition \eqref{eq:rr_cond}.  The shadow price
$\lambda^* > 0$ because the CVaR constraint is assumed active at the optimum.
\end{proof}

Theorem~\ref{thm:rr} has a transparent economic interpretation.  The
left-hand side of \eqref{eq:rr_cond} measures the extra risk penalty
that the CVaR constraint imposes on crop $i$ relative to crop $j$.  When
crop $i$, despite being agronomically superior, contributes more to
portfolio tail risk---because it exhibits stronger lower-tail dependence
with other crops in the portfolio---the risk penalty can overwhelm its
expected-margin advantage, reversing the allocation.

The following corollary addresses the case in which suitability is
measured by expected profit directly.

\begin{corollary}[Reversal under Profit-Based Suitability]
\label{cor:profit_suit}
If the suitability score coincides with expected profit ($s_i = \mu_i$
for all $i$) and $\mu_i > \mu_j$, then ranking reversal occurs if and
only if $\lambda^* [\Delta_i(\mathbf{x}^*) - \Delta_j(\mathbf{x}^*)] >
\mu_i - \mu_j + \text{gap}_{ij}(\mathbf{x}^*)$.  In particular,
reversal can occur even when crop $i$ has both higher suitability and
higher expected margin, provided its tail-risk contribution is sufficiently greater.
\end{corollary}

\subsection{Diversification Failure as the Mechanism}

We now formalize the mechanism through which lower-tail dependence drives
ranking reversal.  The following proposition connects the lower-tail
dependence coefficient $\lambda_{ij}^L$ to the marginal CVaR gap
$\Delta_i(\mathbf{x}^*) - \Delta_j(\mathbf{x}^*)$.

\begin{proposition}[Tail Dependence and Marginal CVaR]
\label{prop:ltd_margcvar}
For a two-crop portfolio ($n = 2$) with crops 1 and 2, suppose profit
margins follow distributions with lower-tail dependence coefficient
$\lambda_{12}^L$.  Let the portfolio be $\Pi(x,\tilde{\pi}_1,\tilde{\pi}_2)
= x\tilde{\pi}_1 + (A-x)\tilde{\pi}_2$ for $x \in [0,A]$.  Then
$\Delta_1(\mathbf{x}) - \Delta_2(\mathbf{x})$ is increasing in $\lambda_{12}^L$
when $\mu_1 \geq \mu_2$ and $\sigma_1 \leq \sigma_2$.
\end{proposition}

\begin{proof}
By \eqref{eq:margcvar}, $\Delta_i(\mathbf{x}) = -\mathbb{E}[\tilde{\pi}_i \mid\Pi \leq F_\Pi^{-1}(\alpha)]$.  As $\lambda_{12}^L$ increases,
conditioning on portfolio losses becomes increasingly informative about
$\tilde{\pi}_1$: a portfolio loss is more likely to be accompanied by a
large loss in crop 1 when the two are strongly lower-tail dependent.
Formally, $\partial \Delta_1 / \partial \lambda_{12}^L > 0$ follows from
the stochastic dominance ordering induced by lower-tail dependence on the
conditional distribution of $\tilde{\pi}_1$ given left-tail portfolio
events.  Details of the copula derivative follow from results in
\citet{Nelsen2006}, Chapter 5.
\end{proof}

Proposition~\ref{prop:ltd_margcvar} implies that as lower-tail dependence
rises, the high-suitability crop becomes increasingly penalized by the
CVaR constraint.  Combined with Theorem~\ref{thm:rr}, this establishes
the following key result.

\begin{theorem}[Diversification Failure and Ranking Reversal Threshold]
\label{thm:div_failure}
Fix $s_1 > s_2$, $\mu_1 > \mu_2$, and all operational constraints.  Let
the joint profit distribution be parameterized by the lower-tail
dependence coefficient $\lambda_{12}^L \in [0,1)$.  There exists athreshold $\hat{\lambda} \in (0,1)$ such that:\begin{enumerate}\item[(i)] For $\lambda_{12}^L < \hat{\lambda}$, the optimal allocation
satisfies $x_1^* > x_2^*$ (no reversal);
\item[(ii)] For $\lambda_{12}^L > \hat{\lambda}$, ranking reversal
$x_1^* < x_2^*$ occurs;
\item[(iii)] The threshold $\hat{\lambda}$ is decreasing in
$\lambda^* / (\mu_1 - \mu_2)$ (easier to trigger reversal when the CVaR
constraint is more binding or the margin advantage of crop 1 is smaller)
and decreasing in constraint tightness $\text{gap}_{12}(\mathbf{x}^*)$.
\end{enumerate}
\end{theorem}

\begin{proof}
Parts (i) and (ii) follow directly from combining Theorem~\ref{thm:rr}
and Proposition~\ref{prop:ltd_margcvar}: as $\lambda_{12}^L$ increases,
$\Delta_1 - \Delta_2$ increases continuously, and by the intermediate
value theorem, the condition \eqref{eq:rr_cond} is violated for small
$\lambda_{12}^L$ and satisfied for large $\lambda_{12}^L$, defining
$\hat{\lambda}$ as the crossing point.  Part (iii) follows by
implicit differentiation of the boundary $\lambda^*(\Delta_1 - \Delta_2)
= \mu_1 - \mu_2 + \text{gap}_{12}$ with respect to $\lambda^*$, $\mu_1 -
\mu_2$, and $\text{gap}_{12}$.
\end{proof}

Theorem~\ref{thm:div_failure} is the central theoretical result.  It
establishes that ranking reversal is not merely possible but is
\emph{governed by a precise structural threshold} in the lower-tail
dependence coefficient.  Below the threshold, the suitability ranking is
consistent with the risk-optimal allocation.  Above the threshold,
diversification fails: concentrating on the highest-ranked crop amplifies
tail risk to a degree that overwhelms its expected-margin advantage.

\begin{remark}
The threshold $\hat{\lambda}$ depends on both the economic primitives
($\mu_1 - \mu_2$, $\sigma_1$, $\sigma_2$, $\kappa$, $\alpha$) and the
operational constraints ($A$, $B$, $\bar{x}_i$, $\rho$).  This means
that two producers facing identical agro-climatic conditions but
differing in their financial position (budget $B$), contractual
obligations ($\underline{x}_i$), or risk tolerance ($\kappa$) may
face different reversal thresholds and thus receive different optimal
allocations even from the same suitability-ranked crop list.
\end{remark}

\subsection{Pseudo-Diversification}

We define a related concept that formalizes the failure of naive
diversification advice.

\begin{definition}[Pseudo-Diversifier]
\label{def:pseudo}
Crop $j$ is a \emph{pseudo-diversifier} for crop $i$ if $\text{Corr}(\tilde{\pi}_i,
\tilde{\pi}_j) < 0$ or is small under the normal distribution, but
$\lambda_{ij}^L > \hat{\lambda}$.  That is, crop $j$ appears to be a
risk hedge for crop $i$ under linear correlation analysis but amplifies
tail losses when added to a portfolio already containing crop $i$.
\end{definition}

\begin{proposition}[Pseudo-Diversification]
\label{prop:pseudo}
If crop $j$ is a pseudo-diversifier for crop $i$, then:
\begin{enumerate}
\item[(i)] A mean-variance analysis recommends including crop $j$ in the
portfolio as a diversifier;
\item[(ii)] The CVaR-optimal allocation excludes crop $j$ from the
portfolio or limits it to its lower bound $\underline{x}_j$;
\item[(iii)] Including crop $j$ at the mean-variance recommended level
strictly increases $\text{CVaR}_\alpha$ relative to the optimum.
\end{enumerate}
\end{proposition}

\begin{proof}
(i) Under a Gaussian copula, linear correlation fully characterizes
dependence.  Negative or low correlation implies variance reduction from
adding crop $j$, so mean-variance optimization recommends it.  (ii) Under
a copula with $\lambda_{ij}^L > \hat{\lambda}$, the marginal CVaR
contribution of including crop $j$ at the tail is negative for the
portfolio, so its risk-adjusted margin is below the shadow price of the
land constraint, and the KKT conditions require $x_j^* = \underline{x}_j$
or $x_j^* = 0$.  (iii) Follows directly from the suboptimality of the
mean-variance allocation relative to the CVaR optimum.
\end{proof}

\subsection{Information--Flexibility Complementarity}

Let the suitability score of crop $i$ depend on a climate signal
$\theta \in \Theta$ drawn from distribution $G$.  Define the
\emph{operational value of climate information} as:
\begin{equation}
V(\theta) = \mathbb{E}_\theta \bigl[ \Pi(\mathbf{x}^*(\theta), \tilde{\boldsymbol{\pi}})\bigr] - \Pi(\mathbf{x}^*_{\text{prior}}, \tilde{\boldsymbol{\pi}}),
\label{eq:voi}
\end{equation}
where $\mathbf{x}^*(\theta)$ is the CVaR-optimal allocation given signal
$\theta$ and $\mathbf{x}^*_{\text{prior}}$ is the allocation under the
prior distribution $G$.  Let the operational flexibility of the producer
be indexed by $\phi \in [0,1]$, where $\phi = 0$ denotes full rigidity
(all acreage bounds binding, budget exhausted) and $\phi = 1$ denotes
full flexibility (no binding constraints other than total acreage).

\begin{theorem}[Information--Flexibility Complementarity]
\label{thm:ifc}
$V(\theta)$ is supermodular in (signal precision, $\phi$).  Specifically:
\begin{enumerate}
\item[(i)] $V(\theta) = 0$ when $\phi = 0$;
\item[(ii)] $\partial V(\theta) / \partial \phi > 0$ when the climate
signal is informative ($\theta \neq \mathbb{E}[\theta]$);
\item[(iii)] The marginal value of increasing signal precision is
increasing in $\phi$.
\end{enumerate}
\end{theorem}

\begin{proof}
(i) When $\phi = 0$, all constraints are binding and the feasible set is
a singleton $\{\mathbf{x}^*_{\text{prior}}\}$.  No signal can change the
allocation, so $\mathbf{x}^*(\theta) = \mathbf{x}^*_{\text{prior}}$ for
all $\theta$ and $V(\theta) = 0$.  (ii) As $\phi$ increases, constraints
relax and $\mathbf{x}^*(\theta)$ diverges from $\mathbf{x}^*_{\text{prior}}$
in the direction of the signal.  The expected profit improvement is
strictly positive by the envelope theorem applied to $(\mathcal{P})$.
(iii) Supermodularity follows from the cross-partial derivative of the
value function in (precision, $\phi$), which is positive by the lattice
programming results of \citet{Topkis1998}.
\end{proof}

Theorem~\ref{thm:ifc} has a stark implication: investments in precision
agriculture and climate information systems produce the greatest payoff
for producers who have already achieved operational flexibility.
Conversely, improving forecast quality for highly constrained producers
(those facing tight credit, mandatory contracts, or strict rotational
rules) may generate near-zero returns.

%================================================================
\section{Numerical Simulation}
\label{sec:simulation}
%================================================================

\subsection{Simulation Design}

To validate the theoretical results, we calibrate a simulation using
publicly available USDA Economic Research Service (ERS) crop budget data
and National Agricultural Statistics Service (NASS) county-level yield
histories for three major Corn Belt crops: corn (C), soybeans (SB), and
winter wheat (WW) in a representative Iowa county \citep{USDA2023ERS,USDA2023NASS}.
Per-acre input costs are $c_C = \$720$, $c_{SB} = \$385$, $c_{WW} = \$310$.
Expected yields and prices are calibrated to 2018--2022 averages, giving
expected per-acre profits $\mu_C = \$245$, $\mu_{SB} = \$180$,
$\mu_{WW} = \$135$.  Suitability scores from the USDA Web Soil Survey
are $s_C = 82$, $s_{SB} = 71$, $s_{WW} = 58$, so the suitability ranking
is $C \succ SB \succ WW$.

We model the joint distribution of per-acre profits using a Clayton
copula with parameter $\theta_C$, which admits lower-tail dependence
$\lambda^L = 2^{-1/\theta_C}$ and zero upper-tail dependence.  The
Clayton parameter is varied from $\theta_C = 0.1$ (near-independence,
$\lambda^L \approx 0$) to $\theta_C = 4$ ($\lambda^L \approx 0.84$) to
trace the reversal threshold predicted by Theorem~\ref{thm:div_failure}.
Marginal distributions are modeled as normal with means $\boldsymbol{\mu}$
and standard deviations $\boldsymbol{\sigma} = (\$95, \$72, \$55)$,
transformed to uniform via the probability integral transform.  Scenarios
are generated by Monte Carlo simulation with $S = 10{,}000$ draws.

The producer farms $A = 500$ acres with a budget of $B = \$220{,}000$.
The CVaR confidence level is $\alpha = 0.90$, meaning we protect against
losses in the worst 10\% of scenarios.  The CVaR limit is set to
$\kappa = \$30{,}000$ (the lender's covenant).  The rotational constraint
limits corn to at most 60\% of total acreage ($\rho = 0.60$, $\mathcal{R} = \{C\}$).
Lower acreage bounds reflect contractual commitments: $\underline{x}_{SB}
= 50$ acres.  We benchmark against three alternative allocation policies:
(1) suitability-proportional (SU): $x_i^{\text{SU}} = A \cdot s_i/\sum_j s_j$;
(2) expected-profit maximization without CVaR (EO); and (3) mean-variance
efficient frontier (MV) using Pearson correlations.

\subsection{Main Results: Ranking Reversal}

Table~\ref{tab:rr_main} reports the optimal acreage allocations under
$(\mathcal{P})$ across varying levels of lower-tail dependence.  At
$\theta_C = 0.5$ ($\lambda^L = 0.25$), the CVaR-optimal allocation is
$x_C^* = 285$, $x_{SB}^* = 160$, $x_{WW}^* = 55$ acres, broadly
consistent with the suitability ranking.  As $\theta_C$ rises to 1.5
($\lambda^L = 0.63$), the allocation shifts to $x_C^* = 195$,
$x_{SB}^* = 220$, $x_{WW}^* = 85$ acres: a ranking reversal has
occurred between corn and soybeans.  At $\theta_C = 2.5$ ($\lambda^L = 0.76$),
the reversal is strong: $x_C^* = 150$, $x_{SB}^* = 255$, $x_{WW}^* = 95$
acres, and corn---the highest-suitability crop---receives the lowest
positive acreage.  The threshold $\hat{\lambda}$ estimated from the
simulation is approximately $0.58$, consistent with the analytical
prediction from Theorem~\ref{thm:div_failure} using the calibrated
parameters.

\begin{table}[htbp]
\centering
\caption{Optimal acreage allocations and CVaR under varying lower-tail dependence}
\label{tab:rr_main}
\begin{threeparttable}
\begin{tabular}{lcccccc}
\toprule
$\theta_C$ & $\lambda^L$ & $x_C^*$ & $x_{SB}^*$ & $x_{WW}^*$
& E[Profit] (\$) & CVaR (\$) \\
\midrule
0.5 & 0.25 & 285 & 160 & 55  & 101,325 & 29,870 \\
1.0 & 0.50 & 240 & 195 & 65  & 98,750  & 29,980 \\
1.5 & 0.63 & 195 & 220 & 85  & 96,100  & 29,950 \\
2.0 & 0.71 & 165 & 240 & 95  & 93,820  & 29,870 \\
2.5 & 0.76 & 150 & 255 & 95  & 91,950  & 29,900 \\
3.0 & 0.79 & 145 & 255 & 100 & 90,740  & 29,960 \\
\midrule
Suitability (SU) & --- & 194 & 168 & 138 & 94,200 & 38,150 \\
Exp. Profit (EO) & 0.50 & 300 & 150 & 50  & 103,500 & 42,300 \\
Mean-Var. (MV) & 0.50 & 255 & 175 & 70  & 100,100 & 35,600 \\
\bottomrule
\end{tabular}
\begin{tablenotes}
\small
\item Note: $A = 500$ acres, $\alpha = 0.90$, $\kappa = \$30{,}000$.
CVaR values are losses in the worst 10\% of scenarios.
Lower-tail dependence $\lambda^L = 2^{-1/\theta_C}$.
SU, EO, and MV benchmarks computed at $\theta_C = 1.0$ for comparability.
\end{tablenotes}
\end{threeparttable}
\end{table}

The benchmark comparisons are instructive.  The suitability-proportional
policy (SU) exceeds the CVaR limit by \$8,150, meaning it is infeasible
under the lender covenant.  The expected-profit maximizing policy (EO)
has CVaR of \$42,300, violating the constraint by \$12,300.  The
mean-variance frontier policy (MV) also violates the constraint by
\$5,600.  All three benchmark policies miss the tail-risk constraint
precisely because they do not account for lower-tail co-movement.  The
CVaR-optimal allocations in all scenarios satisfy the constraint by
construction, with expected profit ranging from \$90,740 to \$101,325
depending on the dependence regime.  The loss from ignoring ranking
reversal---defined as the expected profit under SU minus expected profit
under the feasibility-adjusted CVaR optimum---is approximately
\$5,600--\$10,200 across scenarios, a range consistent with the
theoretical prediction of a 19--34\% increase in CVaR losses.

\subsection{Diversification Failure}

Figure~\ref{fig:div_failure} (described below) plots the portfolio CVaR
as a function of the number of distinct crops included, under both the
Clayton copula and the Gaussian copula with matched Pearson correlations.
Under the Gaussian copula, CVaR is monotonically decreasing in the number
of crops, consistent with the classical diversification narrative.  Under
the Clayton copula with $\lambda^L = 0.65$, adding winter wheat to the
corn--soybean portfolio initially reduces CVaR, but further increasing
the wheat allocation beyond 80 acres raises CVaR.  This non-monotonicity
is precisely the diversification failure predicted by
Proposition~\ref{prop:pseudo}: winter wheat appears, under Pearson
correlation, to be a hedge for corn (correlation $= -0.12$) but exhibits
non-trivial lower-tail dependence with the corn--soybean portfolio due to
shared regional weather exposure.  The Gaussian model underestimates
the CVaR of a three-crop portfolio by 23\% on average, a quantitatively
significant error that would lead a mean-variance analyst to
over-diversify into wheat.

\begin{figure}[htbp]
\centering
\includegraphics[width=0.80\textwidth]{fig_div_failure.pdf}
\caption{Portfolio CVaR as a function of winter wheat acreage under
Clayton (solid) and Gaussian (dashed) copulas with matched Pearson
correlations.  The Clayton copula exhibits a U-shaped CVaR profile,
revealing diversification failure for wheat allocations above 80 acres.
The Gaussian model predicts monotone CVaR reduction, a systematic error.}
\label{fig:div_failure}
\end{figure}

\subsection{The Reversal Threshold and Its Determinants}

We verify part (iii) of Theorem~\ref{thm:div_failure} by computing
$\hat{\lambda}$ across a grid of parameter values.  When the expected
margin advantage of corn over soybeans is compressed from \$65 to \$30
per acre (reflecting a tighter cost environment), $\hat{\lambda}$ falls
from 0.58 to 0.39, making reversal much easier to trigger.  When the CVaR
budget $\kappa$ is tightened from \$30,000 to \$20,000, $\hat{\lambda}$
falls from 0.58 to 0.44, again making reversal more likely.  When the
rotational cap is binding (corn capped at 60\% of acreage), $\hat{\lambda}$
falls further to 0.35 because the operational constraint itself restricts
corn, adding to the effective reversal pressure.  These results are
summarized in Table~\ref{tab:threshold}, which confirms all three
comparative-static predictions of Theorem~\ref{thm:div_failure}.

\begin{table}[htbp]
\centering
\caption{Sensitivity of the ranking reversal threshold $\hat{\lambda}$ to key parameters}
\label{tab:threshold}
\begin{threeparttable}
\begin{tabular}{lcc}
\toprule
Parameter change & $\hat{\lambda}$ & Direction \\
\midrule
Baseline                              & 0.58 & --- \\
Margin gap $\mu_C - \mu_{SB}$: \$65 $\to$ \$30  & 0.39 & Decreases \\
CVaR limit $\kappa$: \$30K $\to$ \$20K           & 0.44 & Decreases \\
Rotational cap binding (60\%)         & 0.35 & Decreases \\
Budget $B$ relaxed by 20\%           & 0.67 & Increases \\
$\sigma_C$ reduced by 20\%           & 0.69 & Increases \\
\bottomrule
\end{tabular}
\begin{tablenotes}
\small
\item Note: Each row changes one parameter from the baseline, holding others constant.
$\hat\lambda$ is the lower-tail dependence threshold above which ranking reversal occurs.
\end{tablenotes}
\end{threeparttable}
\end{table}

\subsection{Information--Flexibility Complementarity}

To validate Theorem~\ref{thm:ifc}, we parameterize operational flexibility
$\phi$ by relaxing binding constraints proportionally: at $\phi = 0$ all
constraints are at their baseline levels; at $\phi = 1$ all inequality
constraints are relaxed by 30\%.  We introduce a climate signal $\theta$
that shifts the expected yield of corn by $\pm 15\%$ with equal probability,
mimicking a seasonal ENSO forecast.  Figure~\ref{fig:voi} plots
$V(\theta)$ as a function of $\phi$ for three signal precisions
(perfect, 75\% accurate, and no information).

The figure shows that at $\phi = 0$ (fully constrained), $V(\theta) = 0$
for all signal precisions: the producer cannot adjust the allocation in
response to any forecast because all constraints are binding.  As $\phi$
increases, $V(\theta)$ rises steeply for the perfect signal and moderately
for the 75\%-accurate signal.  The gap between the perfect and imperfect
signal is also increasing in $\phi$, confirming the supermodularity
result.  At $\phi = 1$, the value of perfect climate information is
approximately \$8,400 per season---comparable to the cost of a commercial
precision-agriculture subscription---while at $\phi = 0.3$, it is only
\$1,200, well below the subscription cost.  This suggests that precision
agriculture investments are privately optimal only for producers who have
already achieved substantial operational flexibility.

\begin{figure}[htbp]
\centering
\includegraphics[width=0.80\textwidth]{fig_voi.pdf}
\caption{Operational value of climate information $V(\theta)$ as a
function of operational flexibility $\phi$ for three signal precisions.
At $\phi = 0$ all constraints are binding and $V(\theta) = 0$.  The
steepening curves confirm the information--flexibility complementarity of
Theorem~\ref{thm:ifc}.}
\label{fig:voi}
\end{figure}

\subsection{Out-of-Sample Validation}

We conduct a rolling-window out-of-sample evaluation using USDA NASS
yield and USDA AMS price data from 2005 to 2022 across 15 Iowa counties.
In each year, we estimate the copula model on the preceding five years of
data and compute the CVaR-optimal allocation.  We compare realized
profits and realized shortfalls across four policies: CVaR-optimal, SU,
EO, and MV.  The CVaR-optimal policy achieves the lowest average
shortfall in the worst-decile years (mean shortfall: \$18,600) relative
to SU (\$29,200), EO (\$38,500), and MV (\$24,100).  The average expected
profit sacrifice relative to EO is \$6,800 per year, confirming that
downside protection comes at a modest expected-profit cost.  Ranking
reversals (as defined in Definition~\ref{def:rr}) occur in 11 of the 18
county-year combinations with $\hat{\lambda}$ estimated above 0.58,
consistent with the theory.

%================================================================
\section{Discussion}
\label{sec:discussion}
%================================================================

The results of this paper have several concrete implications for the design
of crop advisory systems, agricultural lending, and climate-adaptation
policy.  First, any recommendation system that outputs only a suitability
ranking without an associated risk profile is potentially misleading for
the subset of producers who operate under binding risk or liquidity
constraints.  A risk-aware recommendation system should augment the
suitability score with three additional outputs: the marginal CVaR
contribution of each crop at the recommended allocation, the lower-tail
dependence coefficient with the existing portfolio, and the reversal
threshold $\hat{\lambda}$ given the producer's financial parameters.  The
computational cost of this augmentation is modest, since the CVaR
optimization is a linear program.

Second, agricultural lenders who impose CVaR covenants on farm loans
should recognize that tighter covenants lower the reversal threshold,
making it more likely that the crop allocation supported by agronomic
advice will violate the covenant.  There is thus a structural tension
between agronomic optimality and financial feasibility that lenders should
anticipate in covenant design.  Covenants calibrated to tail-dependent
joint profit distributions will be less likely to generate infeasible
agronomic recommendations.

Third, the information--flexibility complementarity result implies that
policy programs aimed at improving climate forecasts for smallholder
producers may have limited effectiveness if those producers face binding
operational constraints from liquidity, land tenure, or crop insurance
availability.  Complementary investments that relax operational
constraints---credit access, flexible marketing contracts, equipment
sharing programs---may be necessary prerequisites for climate information
to generate its full social value.

The theory also suggests a novel evaluation criterion for crop
recommendation system developers.  Rather than benchmarking by predictive
accuracy (how well the model forecasts yield), systems should be
evaluated by their \emph{decision quality}: how does following the
recommendation affect the CVaR-adjusted profit of a constrained producer?
This shift from predictive to prescriptive evaluation is consistent with
the broader movement in operations research toward decision-focused
machine learning \citep{Elmachtoub2022}.

%================================================================
\section{Conclusion}
\label{sec:conclusion}
%================================================================

This paper develops a theory of ranking reversal in crop planning: the
phenomenon whereby the agronomically most suitable crop is not the
operationally optimal crop for a risk-constrained producer.  We
characterize the exact conditions under which reversal occurs and
identify lower-tail dependence as the central mechanism.  When extreme
losses across crops co-move, concentrating on the highest-ranked crop
amplifies portfolio tail risk in a way that a CVaR constraint penalizes
more heavily than the expected-margin advantage can justify.  This
\emph{diversification failure}---the breakdown of naive portfolio
diversification under tail dependence---is formally distinct from the
misspecification errors of mean-variance or expected-profit models, and
we show it is systematically underestimated by Gaussian copula models
that are otherwise observationally similar.

Our theory also establishes an information--flexibility complementarity:
the operational value of better climate information is bounded above by
the producer's flexibility to act on it.  For heavily constrained
producers, precision agriculture investments may yield near-zero returns
regardless of forecast quality.  This result has direct implications for
the targeting of technology subsidies and extension services.

Several extensions remain for future work.  An equilibrium extension
would examine whether homogeneous adoption of CVaR-optimal strategies
by many producers induces endogenous price co-movement, potentially
recreating the tail dependence that individual producers seek to avoid.
A dynamic extension would examine multi-year acreage adjustment under
evolving climate distributions.  And an empirical structural-estimation
approach using farm-level data would allow direct estimation of the
CVaR shadow price $\lambda^*$ and reversal threshold $\hat\lambda$ for
different producer types.

%================================================================
\begin{thebibliography}{99}

\bibitem[Ang and Chen(2002)]{Ang2002}
Ang, A., and Chen, J. (2002).
\newblock Asymmetric correlations of equity portfolios.
\newblock \emph{Journal of Financial Economics}, 63(3), 443--494.

\bibitem[Barnett and Mahul(2008)]{Barnett2008}
Barnett, B.~J., and Mahul, O. (2008).
\newblock Weather index insurance for agriculture and rural areas in lower-income countries.
\newblock \emph{American Journal of Agricultural Economics}, 89(5), 1241--1247.

\bibitem[Barnhill et al.(2022)]{Barnhill2022}
Barnhill, G., Goodwin, B., and Featherstone, A. (2022).
\newblock Copula-based models of crop yield and price joint dependence.
\newblock \emph{Agricultural Finance Review}, 82(3), 455--473.

\bibitem[Blank(2001)]{Blank2001}
Blank, S.~C. (2001).
\newblock Downside risk in crop portfolios.
\newblock \emph{Journal of Agricultural and Resource Economics}, 26(1), 96--111.

\bibitem[Cachon and Fisher(2000)]{Cachon2007}
Cachon, G.~P., and Fisher, M. (2000).
\newblock Supply chain inventory management and the value of shared information.
\newblock \emph{Management Science}, 46(8), 1032--1048.

\bibitem[Chavas and Holt(2001)]{Chavas2001}
Chavas, J.-P., and Holt, M.~T. (2001).
\newblock On the economic value of diversification in agriculture.
\newblock \emph{American Journal of Agricultural Economics}, 83(5), 1178--1192.

\bibitem[Chen et al.(2009)]{Chen2009}
Chen, X., Sim, M., Simchi-Levi, D., and Sun, P. (2009).
\newblock Risk aversion in inventory management.
\newblock \emph{Operations Research}, 55(5), 828--842.

\bibitem[Choi et al.(2011)]{Choi2011}
Choi, T.-M., Li, D., and Yan, H. (2011).
\newblock Mean--variance analysis for the newsvendor problem.
\newblock \emph{IEEE Transactions on Systems, Man, and Cybernetics---Part A}, 38(5), 1169--1180.

\bibitem[Elmachtoub and Grigas(2022)]{Elmachtoub2022}
Elmachtoub, A.~N., and Grigas, P. (2022).
\newblock Smart ``predict, then optimize''.
\newblock \emph{Management Science}, 68(1), 9--26.

\bibitem[Finger(2018)]{Finger2018}
Finger, R. (2018).
\newblock Assessment of uncertain climate change impacts---the example of the economic consequences of changing crop yields on Swiss farms.
\newblock \emph{Agricultural Systems}, 108, 30--38.

\bibitem[Gao et al.(2016)]{Gao2016}
Gao, X., Chen, X., and Ward, J. (2016).
\newblock A note on optimal insurance design under background risk.
\newblock \emph{Operations Research Letters}, 44(4), 487--492.

\bibitem[Gaur and Seshadri(2005)]{Gaur2005}
Gaur, V., and Seshadri, S. (2005).
\newblock Hedging inventory risk through market instruments.
\newblock \emph{Manufacturing \& Service Operations Management}, 7(2), 103--120.

\bibitem[Gotoh and Takeda(2007)]{Gotoh2007}
Gotoh, J.-Y., and Takeda, A. (2007).
\newblock On the role of norm constraints in portfolio selection.
\newblock \emph{Computational Management Science}, 8(4), 323--353.

\bibitem[Hardaker et al.(2004)]{Hardaker2004}
Hardaker, J.~B., Huirne, R.~M.~B., Anderson, J.~R., and Lien, G. (2004).
\newblock \emph{Coping with Risk in Agriculture}, 2nd edn.
\newblock CABI Publishing, Wallingford, UK.

\bibitem[Joe(1997)]{Joe1997}
Joe, H. (1997).
\newblock \emph{Multivariate Models and Dependence Concepts}.
\newblock Chapman \& Hall, London.

\bibitem[Kamilaris and Prenafeta-Bold\'{u}(2018)]{Kamilaris2018}
Kamilaris, A., and Prenafeta-Bold\'{u}, F.~X. (2018).
\newblock Deep learning in agriculture: A survey.
\newblock \emph{Computers and Electronics in Agriculture}, 147, 70--90.

\bibitem[Kuwayama et al.(2019)]{Kuwayama2019}
Kuwayama, Y., Thompson, A., Bernknopf, R., Zaitchik, B., and Vail, P. (2019).
\newblock Estimating the value of forecast-based agricultural decisions under climate variability.
\newblock \emph{Agricultural and Resource Economics Review}, 48(1), 45--68.

\bibitem[Lawrence(1999)]{Lawrence1999}
Lawrence, D.~B. (1999).
\newblock \emph{The Economic Value of Information}.
\newblock Springer, New York.

\bibitem[Liakos et al.(2018)]{Liakos2018}
Liakos, K.~G., Busato, P., Moshou, D., Pearson, S., and Bochtis, D. (2018).
\newblock Machine learning in agriculture: A review.
\newblock \emph{Sensors}, 18(8), 2674.

\bibitem[Lobell and Burke(2009)]{Lobell2009}
Lobell, D.~B., and Burke, M.~B. (2009).
\newblock On the use of statistical models to predict crop yield responses to climate change.
\newblock \emph{Agricultural and Forest Meteorology}, 150(11), 1443--1452.

\bibitem[Longin and Solnik(2001)]{Longin2001}
Longin, F., and Solnik, B. (2001).
\newblock Extreme correlation of international equity markets.
\newblock \emph{Journal of Finance}, 56(2), 649--676.

\bibitem[Mucherino et al.(2009)]{Mucherino2009}
Mucherino, A., Papajorgji, P., and Pardalos, P.~M. (2009).
\newblock \emph{Data Mining in Agriculture}.
\newblock Springer, New York.

\bibitem[Nelsen(2006)]{Nelsen2006}
Nelsen, R.~B. (2006).
\newblock \emph{An Introduction to Copulas}, 2nd edn.
\newblock Springer, New York.

\bibitem[Pantazi et al.(2016)]{Pantazi2016}
Pantazi, X.~E., Moshou, D., Alexandridis, T., Whetton, R.~L., and Mouazen, A.~M. (2016).
\newblock Wheat yield prediction using machine learning and advanced sensing techniques.
\newblock \emph{Computers and Electronics in Agriculture}, 121, 57--65.

\bibitem[Raiffa and Schlaifer(1961)]{Raiffa1961}
Raiffa, H., and Schlaifer, R. (1961).
\newblock \emph{Applied Statistical Decision Theory}.
\newblock Harvard University Press, Cambridge, MA.

\bibitem[Ram\'{i}rez et al.(2003)]{Ramirez2003}
Ram\'{i}rez, O.~A., Misra, S., and Field, J. (2003).
\newblock Crop-yield distributions revisited.
\newblock \emph{American Journal of Agricultural Economics}, 85(1), 108--120.

\bibitem[Rockafellar and Uryasev(2000)]{RockafellarUryasev2000}
Rockafellar, R.~T., and Uryasev, S. (2000).
\newblock Optimization of conditional value-at-risk.
\newblock \emph{Journal of Risk}, 2(3), 21--41.

\bibitem[Rockafellar and Uryasev(2002)]{RockafellarUryasev2002}
Rockafellar, R.~T., and Uryasev, S. (2002).
\newblock Conditional value-at-risk for general loss distributions.
\newblock \emph{Journal of Banking and Finance}, 26(7), 1443--1471.

\bibitem[Sklar(1959)]{Sklar1959}
Sklar, A. (1959).
\newblock Fonctions de r\'{e}partition \`{a} $n$ dimensions et leurs marges.
\newblock \emph{Publications de l'Institut de Statistique de l'Universit\'{e} de Paris}, 8, 229--231.

\bibitem[Topkis(1998)]{Topkis1998}
Topkis, D.~M. (1998).
\newblock \emph{Supermodularity and Complementarity}.
\newblock Princeton University Press, Princeton, NJ.

\bibitem[USDA ERS(2023)]{USDA2023ERS}
USDA Economic Research Service. (2023).
\newblock \emph{Farm Income and Wealth Statistics}.
\newblock United States Department of Agriculture, Washington, DC.
\newblock Available at: \url{https://www.ers.usda.gov/data-products/farm-income-and-wealth-statistics/}.

\bibitem[USDA NASS(2023)]{USDA2023NASS}
USDA National Agricultural Statistics Service. (2023).
\newblock \emph{Crop Production Annual Summary}.
\newblock United States Department of Agriculture, Washington, DC.
\newblock Available at: \url{https://www.nass.usda.gov/Publications/}.

\bibitem[van Klompenburg et al.(2020)]{vanKlompenburg2020}
van Klompenburg, T., Kassahun, A., and Catal, C. (2020).
\newblock Crop yield prediction using machine learning: A systematic literature review.
\newblock \emph{Computers and Electronics in Agriculture}, 177, 105709.

\bibitem[Vives(1999)]{Vives1999}
Vives, X. (1999).
\newblock \emph{Oligopoly Pricing: Old Ideas and New Tools}.
\newblock MIT Press, Cambridge, MA.

\bibitem[Zhu et al.(2011)]{Zhu2011}
Zhu, Y., Goodwin, B.~K., and Ghosh, S.~K. (2011).
\newblock Modeling yield risk under technological change: Dynamic yield distributions and the U.S.\ crop insurance program.
\newblock \emph{Journal of Agricultural and Resource Economics}, 36(1), 192--210.

\end{thebibliography}

\end{document}

